\chapter*{ประวัติผู้เขียน}
\addcontentsline{toc}{chapter}{ประวัติผู้เขียน}

\vspace{0.4cm}

% ส่วนบน: Profile Header ระดับพรีเมียม
\noindent
\begin{minipage}[c]{0.22\textwidth}
  \centering
  \begin{tikzpicture}
    % ขอบทองราชภัฏพร้อมเงา
    \draw[draw=rbruGold, line width=2pt, fill=rbruNavy!15] (0, 0) circle (1.5cm);
    \clip (0, 0) circle (1.46cm);
    % ไอคอนภาพจำลองนักวิจัยฟิสิกส์และเทคโนโลยี
    \node at (0, 0) {
      \includegraphics[width=3.0cm, height=3.0cm, keepaspectratio]{figures/app_icon.png}
    };
  \end{tikzpicture}
\end{minipage}
\hfill
\begin{minipage}[c]{0.75\textwidth}
  {\fontsize{18pt}{22pt}\selectfont\bfseries\color{rbruNavy} ผู้ช่วยศาสตราจารย์ ดร.ชีวะ ทัศนา}\\[2mm]
  {\fontsize{11pt}{15pt}\selectfont\bfseries\color{rbruSlate} หน่วยวิเคราะห์-ติดตามคุณภาพดินและสภาพอากาศเชิงพื้นที่}\\[1mm]
  {\fontsize{10.5pt}{14.5pt}\selectfont คณะวิทยาศาสตร์และเทคโนโลยี มหาวิทยาลัยราชภัฏรำไพพรรณี}\\[2mm]
  {\fontsize{9.5pt}{13pt}\selectfont\color{rbruCyan!90!black} \textbf{Email} chewa.t@rbru.ac.th \quad\textbar\quad \textbf{Google Scholar ID} \texttt{jJz3B6AAAAAJ}}
\end{minipage}

\vspace{0.6cm}
{\color{rbruSlate!30}\hrule height 1pt}
\vspace{0.6cm}

% ส่วนล่าง: Balanced Two-Column Body พอดี 1 หน้าเดียว
\noindent
\begin{minipage}[t]{0.48\textwidth}
  {\fontsize{13pt}{16pt}\selectfont\bfseries\color{rbruNavy} ประวัติการศึกษา}\\[2mm]
  \begin{itemize}[leftmargin=1.5em, itemsep=2pt, topsep=1pt]
    \item \textbf{ปร.ด. (ฟิสิกส์ประยุกต์)}\\
          สถาบันเทคโนโลยีพระจอมเกล้าเจ้าคุณทหารลาดกระบัง (KMITL)
    \item \textbf{วท.ม. (ฟิสิกส์ประยุกต์)}\\
          สถาบันเทคโนโลยีพระจอมเกล้าเจ้าคุณทหารลาดกระบัง (KMITL)
    \item \textbf{วท.บ. ฟิสิกส์ (คอมพิวเตอร์และอิเล็กทรอนิกส์)}\\
          มหาวิทยาลัยนเรศวร
  \end{itemize}

  \vspace{4mm}
  {\fontsize{13pt}{16pt}\selectfont\bfseries\color{rbruNavy} รายวิชาที่รับผิดชอบสอน}\\[2mm]
  \begin{itemize}[leftmargin=1.5em, itemsep=2pt, topsep=1pt]
    \item ฟิสิกส์สำหรับวิทยาศาสตร์สุขภาพและกีฬา
    \item การประมวลผลสัญญาณและเซนเซอร์อัจฉริยะ
    \item ปัญญาประดิษฐ์และการเรียนรู้ของเครื่อง
    \item การพัฒนาแอปพลิเคชันบนอุปกรณ์พกพา
  \end{itemize}
\end{minipage}
\hfill
\begin{minipage}[t]{0.48\textwidth}
  {\fontsize{13pt}{16pt}\selectfont\bfseries\color{rbruNavy} ความเชี่ยวชาญและงานวิจัยที่สนใจ}\\[2mm]
  \begin{itemize}[leftmargin=1.5em, itemsep=2pt, topsep=1pt]
    \item \textbf{การประมวลผลสัญญาณชีวการแพทย์และจลนศาสตร์}\\
          (Biomedical Signal Processing and Human Kinematics)
    \item \textbf{เซนเซอร์ตรวจวัดความเร่งและอินเทอร์เน็ตของสรรพสิ่ง}\\
          (MEMS Accelerometers, IMU and Edge Health IoT)
    \item \textbf{ชีวกลศาสตร์และวิทยาศาสตร์การออกกำลังกายประยุกต์}\\
          (Applied Biomechanics and Exercise Science)
    \item \textbf{สถาปัตยกรรมส่วนต่อประสานผู้ใช้ทางวิทยาศาสตร์}\\
          (Scientific Responsive UI/UX and Bento Grid Design)
    \item \textbf{ฟิสิกส์เกษตรดิจิทัลและปัญญาประดิษฐ์เชิงพื้นที่}\\
          (Digital Agriphysics and Precision Spatial AI)
  \end{itemize}
\end{minipage}

\vfill
\begin{center}
  {\fontsize{9pt}{12pt}\selectfont\color{gray} เอกสารคู่มือวิชาการและวิจัยฉบับสมบูรณ์ มหาวิทยาลัยราชภัฏรำไพพรรณี พุทธศักราช 2569}
\end{center}
